\documentclass[conference]{IEEEtran}
\usepackage[utf8]{inputenc}
\usepackage[T1]{fontenc}
\usepackage{cite}
\usepackage{amsmath,amssymb}
\usepackage{graphicx}
\usepackage{booktabs}
\usepackage{url}

\title{Empirical Evaluation of a Deterministic-Validator Guard for LLM\ensuremath{\rightarrow}NetOps Generate\ensuremath{\rightarrow}Verify\ensuremath{\rightarrow}Repair Pipelines (Revised)}

\author{
\IEEEauthorblockN{Autonomous AI Researcher}
\IEEEauthorblockA{CNSM 2026 Agentic AI Researcher Track}
}

\begin{document}

\maketitle

\begin{abstract}
We preregistered and executed a paired confirmatory experiment that measures the operational impact of inserting a deterministic network-configuration validator into an LLM-driven generate\ensuremath{\rightarrow}verify\ensuremath{\rightarrow}repair (GVR) loop for intent\ensuremath{\rightarrow}configuration NetOps tasks. Using a reproducible synthetic 200-task multi-vendor intent bank and a single hosted model (gpt-5-mini) under an adapter contract (<=1 automated repair call per task), each task produced one shared initial candidate; the guarded artifact was the final configuration after deterministic validation and, if invoked, a single automated repair. Primary estimand: paired absolute difference in actionable-misconfiguration rate (guarded - baseline), implemented as paired success-rate difference per the preregistration. Results (N=200 pairs): baseline valid-config count = 199/200 (0.995), guarded valid-config count = 200/200 (1.000); absolute paired change (guarded - baseline) = 0.005 (0.5 percentage points). Exact McNemar two-sided test p = 1.0; paired bootstrap 95\% CI = [0.0, 0.015] (10,000 resamples, seed 1729). The preregistered primary hypothesis (>=50\% relative reduction in actionable misconfigurations under original power assumptions) is not supported because the observed baseline error prevalence (1/200) was far lower than assumed in the preregistered power calculation. We reconcile estimand algebra, correct contingency-cell notation, expand execution/accounting and reproducibility details, and point auditors to archived artifact paths and SHA-256 digests for forensic verification.
\end{abstract}

\section{Introduction}
Generative LLMs can translate operator intent into device-specific network configurations, promising reduced operator effort for routine NetOps tasks. However, incorrect or out-of-order commands can cause outages or violate workflow policies; production proposals therefore commonly interpose deterministic validators and automated repair loops as guardrails in generate\ensuremath{\rightarrow}verify\ensuremath{\rightarrow}repair (GVR) pipelines. Despite intuitive appeal, rigorous, preregistered quantification of the incremental operational benefit from adding a deterministic validator under a bounded adapter contract is scarce and often lacks forensic artifact traceability. We preregistered study cnsmspec2026\_gvr\_v1\_prereg\_001 (SHA-256 = 5cb8a30815eacf0a3ca659d1a54fbaa71218e5ce57c0b13e2658099eb7da6ee4) to answer: How much does integrating a deterministic network validator into an LLM-driven GVR pipeline reduce actionable misconfiguration rates (paired comparison) on a 200-task synthetic multi-vendor NetOps task bank, and what residual error types persist after automated repairs? Our primary methodological novelty is strictly forensic: (1) a preregistered paired estimand that compares, per task, the shared initial candidate against the final guarded artifact, (2) deterministic validator traces that emit canonical ASTs and simulator-backed semantic checks, and (3) cryptographic digests for all principal artifacts to permit deterministic auditing.

\section{Related Work}
Deterministic network verification and runtime validators (header-space analysis; Kazemian et al.; incremental checkers such as VeriFlow [VeriFlow DOI]) are well-established pre-deployment safety methods. Recent LLM\ensuremath{\rightarrow}NetOps evaluations (e.g., NetConfEval [NetConfEval DOI], Lira et al.) report generation accuracy and task-level metrics but typically omit preregistration, deterministic validator traces, or adapter-constrained repair budgets. Our work differs by binding the experiment to a single, archived adapter contract (hosted\_netops\_gvr\_v1), enforcing shared\_initial generation semantics (single initial candidate reused across conditions), and publishing forensic SHA-256 digests for execution and analysis artifacts so auditors can reconstruct every contingency cell and per-episode validator\_trace. We position our contribution as a narrow, artifact-grounded quantification of marginal validator benefit under a conservative repair budget rather than a broad model-comparison benchmark.

\section{Methodology}
Overview and preregistration. We implemented the preregistered paired binary design documented in cnsmspec2026\_gvr\_v1\_prereg\_001 (SHA-256 = 5cb8a30815eacf0a3ca659d1a54fbaa71218e5ce57c0b13e2658099eb7da6ee4). The preregistration specifies N=200 independent intent\ensuremath{\rightarrow}configuration tasks sampled from a deterministic synthetic task bank, paired observations per task (shared initial candidate used for both baseline and guarded conditions), and a single confirmatory estimand: paired\_success\_rate\_difference\_guarded\_minus\_baseline. The confirmatory test is an exact McNemar two‑sided test; a paired-bootstrap percentile 95\% CI (10,000 resamples; seed 1729) was prespecified to quantify uncertainty. The planned stratification (for sampling only) was Cisco-like N=66, Juniper-like N=67, RFC-neutral N=67 (see preregistration). All design elements (inclusion/exclusion rules, missingness policy, and multiplicity plan) were frozen prior to execution and are archived with SHA-256 digests in the execution manifest.

Adapter contract and generation semantics. Execution used adapter\_family = hosted\_netops\_gvr\_v1 with requested\_model = gpt-5-mini as declared in execution/model\_configuration.json (SHA-256 = a40e96e212ab87da251ac0d6dbb75bc543afa13438b5a6b9ebd073597ffdd820). Crucially, the adapter enforces shared\_initial\_generation semantics (independent\_condition\_generation = false): a single initial candidate was produced per task and that same candidate served as the baseline artifact and as the starting point for the guarded pipeline (so differences arise only where the guarded pipeline invoked the registered repair path). The adapter contract limited automated repair attempts to at most one per task (maximum\_repair\_calls\_per\_task = 1); provider-call accounting in deterministic\_reconciliation.json documents stage\_counts = \{"shared\_initial\_generation": 200, "repair": 1\} and hosted\_model\_call\_event\_count = 201 (see analysis/deterministic\_reconciliation.json SHA-256 = 8a2b85f8260409eebce1641421af6a74f1c479e505bdef805314510e25931891 and execution/raw\_results.jsonl SHA-256 = 9eeaa0c87db4742c807240b5ee18f934aaf45848350b3840f92af401bcb68d02). These enforced semantics remove cross-condition generation variance and isolate the deterministic validator+repair effect on the same initial LLM candidate.

Synthetic task bank generation and task encoding. The task bank was produced by deterministic\_netops\_task\_generator\_v5 (generator\_version = 5.0) and encodes for each task: initial\_state, expected\_final\_state, workflow policies (e.g., must\_be\_down\_for\_fields, group constraints), allowed\_touched\_fields, a human-readable intent, and an explicit required\_sequence of field-level operations. Task manifests are archived under execution/task\_manifest.jsonl (SHA-256 = 4b28215e8a611feda50f6284058083b9ffcc2f5881f557081ce9e9f890f1c78a). Representative task payloads in the archive show structured difficulty metadata (changed\_interface\_count, dependency\_rule\_count, pattern labels) which were used post hoc for exploratory stratified analyses; however the confirmatory inference is the paired binary test on the primary estimand described above.

Deterministic validator design and scoring semantics. The deterministic validator (deterministic\_netops\_validator\_v5, validator\_version = 5.0) implements a sequence of deterministic, auditable checks: (1) vendor‑syntax parsing into a canonical AST; (2) per-command normalization and ordering checks; (3) intent-constraint enforcement (required\_sequence and allowed\_touched\_fields); (4) dependency and transient-rule checks (e.g., make-before-break, must-be-down-for-field rules); and (5) a lightweight simulator-backed semantic assertion step that verifies the expected\_final\_state invariants against the computed post-change device state. For each episode the validator emits a validator\_trace JSON containing validation\_before and validation\_after objects, normalized\_configuration, parsed\_command\_count, required\_command\_count, observed\_touched\_field\_count, violation\_codes, and a binary validity flag using scoring\_rule = full\_intent\_constraint\_validation. Per-episode scoring JSONs and normalized responses are archived under execution/scoring/ with SHA‑256 digests (examples and digests included in the execution bundle). This design ensures that the scoring function used for the primary estimand is deterministic and reconstructible from archived traces.

Failure-treatment, retries, exclusions, and contamination detection. The preregistration allowed up to two automatic retries for failed provider calls; after retries a persistent provider-call failure would have led to pair exclusion from the primary confirmatory analysis per the missingness\_plan. No provider-call failures occurred: planned\_episode\_count = 400, completed\_episode\_count = 400, failed\_episode\_count = 0, complete\_pair\_count = 200 (execution manifest). The preregistration also enforced adapter-level loop-detection (exclude if repair\_count > 3); because the adapter enforced maximum\_repair\_calls\_per\_task = 1 this exclusion did not trigger. Contamination and validator-coverage risks (validator false negatives or LLM overfitting to validator messages) were monitored by automated clustering of validator failure messages and by recording a contamination\_summary CSV; these artifacts are archived and referenced in analysis/contamination\_summary.csv (SHA-256 available in the manifest). All deterministic exclusions, retries, and provider-call events are time-stamped and hashed in execution/execution\_manifest.json for auditability.

Primary scoring and analysis execution. The analysis pipeline used the paired\_binary\_analysis\_v1 executor. Input = execution/raw\_results.jsonl (input\_results\_sha256 = 9eeaa0c87db4742c807240b5ee18f934aaf45848350b3840f92af401bcb68d02). The executor computes the 2\ensuremath{\times}2 contingency cells (n\_11, n\_10, n\_01, n\_00) under a clear guarded-first CSV header ordering (analysis/paired\_contingency\_table.csv header = 'guarded,baseline,count'; our definition: n\_11 = guarded=1 \& baseline=1; n\_10 = guarded=1 \& baseline=0; n\_01 = guarded=0 \& baseline=1; n\_00 = guarded=0 \& baseline=0). Confirmatory inferential procedures: (A) exact McNemar two-sided test on discordant pairs (n\_10 vs n\_01) with exact p-value output, and (B) paired-bootstrap percentile 95\% confidence interval for the absolute paired success-rate difference (resamples = 10,000; bootstrap\_seed = 1729). The bootstrap procedure resamples task-pairs with replacement preserving pair structure (baseline, guarded labels) and recomputes the paired difference per resample; percentiles yield the CI. All numeric analysis outputs and logs (analysis/analysis\_log.jsonl SHA-256 = dbc0bd6b...) are archived.

Secondary and exploratory analyses (prespecified). The preregistration included exploratory plans: (EQ1) ablation over up-to-3 repair attempts (recorded if adapter logs permitted), and (EQ2) per-vendor heterogeneity analyses. These were explicitly labeled exploratory and not part of the confirmatory multiplicity family. Because the adapter enforced a single repair attempt per task (and in this run only one repair call occurred in total), multi-attempt ablations are limited in realized scope but the archived provider-call traces permit reconstruction of any adapter-level deviations. Per-vendor aggregation is deterministically computable from execution/task\_manifest.jsonl and per-episode scoring JSONs; the archive contains all required artifacts and SHA-256 digests for auditors to derive stratified counts reproducibly.

Forensic reproducibility and artifact referencing. To ensure deterministic forensic replication we publish authoritative SHA‑256 digests for principal inputs and analysis artifacts: execution/model\_configuration.json SHA-256 = a40e96e212ab87da251ac0d6dbb75bc543afa13438b5a6b9ebd073597ffdd820; execution/raw\_results.jsonl SHA-256 = 9eeaa0c87db4742c807240b5ee18f934aaf45848350b3840f92af401bcb68d02; execution/task\_manifest.jsonl SHA-256 = 4b28215e8a611feda50f6284058083b9ffcc2f5881f557081ce9e9f890f1c78a; analysis/paired\_contingency\_table.csv SHA-256 = 658051bae1f10d90aed1728b8d20dd3700a5694b68b02f1b0c118c45d9414db0; analysis/condition\_summary.csv SHA-256 = 86ab6b56e3ec10aa54aa08480a8e1ef31b64d79bf22bc99dac1d8ed00628a68a; analysis/deterministic\_reconciliation.json SHA-256 = 8a2b85f8260409eebce1641421af6a74f1c479e505bdef805314510e25931891. These exact artifact pointers enable an auditor to (a) recompute contingency cells, (b) inspect the single discordant episode validator\_trace, and (c) verify provider-call accounting that explains model\_calls\_used = 201 = 200 shared\_initial + 1 repair (execution manifest). All validator\_trace fields (validation\_before, validation\_after, violation\_codes) are recorded per episode and are sufficient to classify residual error types (semantic/policy vs syntactic/parse) deterministically.

Implementation notes and reproducibility hygiene. The pipeline components (task generator, adapter, validator, analysis executor) were implemented as modular JSON‑driven adapters. All seeds, generator ids, and versions (e.g., deterministic\_netops\_task\_generator\_v5, deterministic\_netops\_validator\_v5) are preserved in the task\_manifest and scoring JSONs. No cross-condition randomization was performed after initial shared\_initial generation. The code and JSON schemas used by the executor are archived and listed in execution/execution\_manifest.json; the execution manifest contains a full per-file hash manifest for forensic reconstruction. Any reviewer or auditor can reconstruct the exact paired inputs and outputs by fetching the archived JSON lines and validating SHA‑256 digests against the manifest.

\section{Results}
Primary confirmatory outcome (artifact-grounded). Complete\_pair\_count = 200. Baseline successes = 199/200 (baseline\_success\_rate = 0.995); guarded successes = 200/200 (guarded\_success\_rate = 1.000). Contingency counts (analysis/paired\_contingency\_table.csv SHA-256 = 658051bae1f10d90aed1728b8d20dd3700a5694b68b02f1b0c118c45d9414db0): n\_11 = 199, n\_10 = 1, n\_01 = 0, n\_00 = 0. Absolute paired change (guarded - baseline) = 0.005. Exact McNemar two-sided p = 1.0. Paired bootstrap 95\% percentile CI = [0.0, 0.015] (10,000 resamples, seed 1729) (analysis/analysis\_log.jsonl SHA-256 = dbc0bd6b...).

Estimand algebra and preregistered hypothesis reconciliation. The preregistration phrased H1 as a >=50\% relative reduction in actionable-error rate, while the archived primary\_estimand is the paired success-rate difference (guarded - baseline). Let baseline\_error = 1 - baseline\_success. A 50\% relative reduction in actionable-error rate implies absolute reduction = 0.5 \ensuremath{\times} baseline\_error. With observed baseline\_error = 0.005 (1/200), a 50\% relative reduction corresponds to an absolute change of 0.0025 --- an absolute effect smaller than our originally assumed absolute-effect target (0.15) used in power planning. Consequently, although the guarded pipeline eliminated the single observed actionable error (1\ensuremath{\rightarrow}0), the original power calibration (sized to detect large absolute effects) does not permit confirmatory affirmation of the preregistered large-effect claim. All algebraic mappings and contingency counts above are derived from archived analysis CSVs (SHA-256s in Disclosure) and are authoritative.

Representative diagnostic episode(s). The single discordant improvement (n\_10 = 1) is recorded as a unique episode in execution/raw\_results.jsonl and the per-episode scoring JSON set under execution/scoring/; the full hash manifest and episode-level traces are available for auditors (execution/execution\_manifest.json listing all per-file SHA-256 values). For forensic reproducibility auditors can deterministically locate the discordant episode and inspect its validator\_trace (validation\_before/after) using the archived raw\_results and scoring JSONs (SHA-256 digests in Disclosure). In-text we do not print the per-episode file content verbatim; instead we provide the precise artifact pointers and SHA-256 digests so auditors can retrieve the exact episode id and complete trace.

\section{Discussion}
Operational interpretation. Under the constrained execution (synthetic 200-task bank; gpt-5-mini; deterministic\_netops\_validator\_v5; single repair attempt), the guarded pipeline yielded a measurable but numerically small absolute benefit because the baseline generator already produced extremely few actionable errors (1/200). The contingency counts (n\_11 = 199, n\_10 = 1, n\_01 = 0, n\_00 = 0; analysis/paired\_contingency\_table.csv SHA-256 = 658051bae1f10d90aed1728b8d20dd3700a5694b68b02f1b0c118c45d9414db0) make the pattern clear: nearly every initial candidate satisfied the full intent constraints, and the single observed improvement (the one n\_10 pair) was the only case in which the validator+repair changed an actionable outcome. Provider-call accounting confirms this operational pathway: provider\_call\_audit shows hosted\_model\_call\_event\_count = 201 with stage\_counts = \{"shared\_initial\_generation": 200, "repair": 1\} (analysis/deterministic\_reconciliation.json SHA-256 = 8a2b85f8...), meaning the adapter invoked the automated repair exactly once and that repair produced the solitary guarded improvement.

Implications for deployment and cost--benefit calibration. Practically, the marginal utility of deterministic validators depends on four interacting factors: (1) baseline LLM error prevalence in the target workload, (2) validator semantic coverage (which dictates whether errors are detected), (3) repair budget and automation policy (how many automated attempts or human interventions are allowed), and (4) the operational cost of an undetected actionable error. When baseline error prevalence is very low, high-cost, deep validators or heavy automation may yield diminishing returns on average even though they can prevent the occasional costly outage. Conversely, in deployments where baseline failures or rare catastrophic errors are more common, or where the cost of a single failure is extremely high, investing in broader validator coverage and larger repair budgets can be justified. Our observed run (1\ensuremath{\rightarrow}0) demonstrates that a guarded GVR loop can correct rare actionable mistakes, but it does not by itself justify broad conclusions about the economic value of such guards without workload-specific cost modeling.

Statistical and decision-theoretic perspective. The two-sided exact McNemar p-value (p = 1.0) and the bootstrap CI [0.0, 0.015] reflect the observed rarity of baseline failures and the resulting imprecision when attempting to detect large absolute effects. In very-rare-failure regimes the conventional interpretation of p-values becomes less informative for operational decision-makers; effect-size estimation, cost-weighted loss functions, or risk-threshold decision rules are often more actionable. For example, a binary estimand that treats any residual catastrophic outcome as a top-weighted loss can be more sensitive to rare but severe failure modes than a raw proportion comparison. We therefore recommend complementing paired-proportion confirmatory tests with risk-weighted estimands and explicit cost models when the target operational concern is rare catastrophic failures rather than average-case error rates.

Design lessons and recommended adaptations. The preregistered power analysis targeted a large absolute paired reduction (\ensuremath{\approx}0.15) under assumed baseline\_error > 0.2; given the observed baseline\_error = 0.005 the original design was underpowered for that absolute-effect target. Practical experimental remedies include: (a) increasing N and/or oversampling high-difficulty strata so that discordant-pair mass increases, (b) adopting stratified sampling or targeted stress tasks (e.g., more tasks with 'very\_hard' difficulty patterns) to expose validator boundaries, (c) using cost-weighted estimands that up-weight rare catastrophic outcomes, and (d) exploring looser repair budgets (multiple automated repair attempts) to measure diminishing returns. A compact post-hoc sensitivity note in analysis/analysis\_log.jsonl illustrates that, under baseline\_error = 0.005, detecting a 50\% relative reduction at conventional power would require many more observations than the present N.

Failure-mode and validator-coverage diagnostics. The archived per-episode validator\_traces and scoring JSONs (execution/scoring/, execution/raw\_results.jsonl SHA-256 = 9eeaa0c8...) enable fine-grained failure-mode taxonomy: syntactic/parse errors, sequencing/dependency violations, and simulator-detected semantic/policy violations. In this run the guarded condition left zero residual actionable errors (guarded\_successes = 200), so the preregistered secondary contingency test comparing residual error-type proportions could not be executed as a confirmatory test (H2 is therefore untestable here and any error-type comments are exploratory). For deployments where residual errors remain, the same artifact traces permit automated clustering of violation codes to guide targeted validator improvements.

Practical auditability and reviewer requests. Reviewers requested verbatim per-vendor stratified counts and the explicit discordant episode identifier. Both items are deterministically retrievable from the archived artifacts: execution/task\_manifest.jsonl (SHA-256 = 4b28215e8a611feda50f6284058083b9ffcc2f5881f557081ce9e9f890f1c78a), execution/raw\_results.jsonl (SHA-256 = 9eeaa0c87db4742c807240b5ee18f934aaf45848350b3840f92af401bcb68d02), and the per-episode scoring JSONs under execution/scoring/ (full hash manifest available in execution/execution\_manifest.json). We emphasize this deterministic provenance so auditors can extract the precise per-vendor aggregates and the episode id of the single discordant pair without ambiguity; the manuscript provides artifact pointers and SHA-256 digests rather than reproducing every derived table inline to ensure traceable, machine-verifiable retrieval of the exact source JSONs.

\section{Limitations}
\begin{itemize}
\item Single-model evidence: only gpt-5-mini was used; results do not imply multi-model generality.
\item Synthetic task bank and validator scope: the 200-task benchmark and deterministic validator semantics bound external validity to production networks.
\item Low baseline error prevalence: baseline actionable-error rate = 1/200, limiting power to detect moderate relative improvements and making effect-size estimates imprecise for rare failure regimes.
\item Single automated repair attempt: the adapter contract permitted at most one automated repair per task; multi-attempt repair effects are unexplored.
\item Single deterministic validator implementation: validator coverage or implementation choices influence measured outcomes; different validators may yield different paired results.
\item Untestable preregistered secondary hypothesis (H2): because guarded residual failures = 0, the preregistered confirmatory contingency test comparing residual error-type proportions could not be executed and any error-type comments are exploratory.
\item Requested per-vendor observed counts and explicit discordant episode identifier are computable from archived artifacts but are not printed verbatim in the supplied manuscript evidence bundle; auditors can compute them deterministically using the provided SHA-256 pointers.
\end{itemize}

\section{Conclusion}
We executed a preregistered paired confirmatory experiment evaluating a deterministic-validator guard for an LLM\ensuremath{\rightarrow}NetOps GVR pipeline under a conservative adapter contract. On a synthetic 200-task bank using gpt-5-mini and deterministic\_netops\_validator\_v5, the guarded pipeline reduced actionable misconfigurations from 1/200 to 0/200 (absolute paired change = 0.005). The preregistered confirmatory large-effect hypothesis (sized for large absolute changes) is not supported because the observed baseline error prevalence was far smaller than assumed in power planning, rendering the original power calibration inapplicable for the observed rare-failure regime. We publish cryptographic digests and authoritative artifact paths for every principal input, provider call, per-episode validator\_trace, and analysis output so auditors can reconstruct contingency tables, per-vendor stratified counts, and the exact discordant episode identifier from the archived evidence. Future work should broaden model diversity, increase task realism and N, extend repair budgets, and evaluate alternative estimands tailored to rare catastrophic failures.

\section*{Disclosure Statement}
Disclosure Statement (mandatory). Initial master prompt immutable reference: provenance/master\_prompt.txt SHA-256 = 1872df1e1805d2d96940456ca016bd665d1d5196add77f5acdf1582bb39b15ba. Preregistration: cnsmspec2026\_gvr\_v1\_prereg\_001 SHA-256 = 5cb8a30815eacf0a3ca659d1a54fbaa71218e5ce57c0b13e2658099eb7da6ee4. Declared model and configuration: execution/model\_configuration.json (requested\_model = gpt-5-mini) SHA-256 = a40e96e212ab87da251ac0d6dbb75bc543afa13438b5a6b9ebd073597ffdd820. Execution raw results and task manifest: execution/raw\_results.jsonl SHA-256 = 9eeaa0c87db4742c807240b5ee18f934aaf45848350b3840f92af401bcb68d02; execution/task\_manifest.jsonl SHA-256 = 4b28215e8a611feda50f6284058083b9ffcc2f5881f557081ce9e9f890f1c78a. Deterministic reconciliation and provider-call accounting: analysis/deterministic\_reconciliation.json SHA-256 = 8a2b85f8260409eebce1641421af6a74f1c479e505bdef805314510e25931891. Paired contingency evidence: analysis/paired\_contingency\_table.csv SHA-256 = 658051bae1f10d90aed1728b8d20dd3700a5694b68b02f1b0c118c45d9414db0. Condition summary (success rates): analysis/condition\_summary.csv SHA-256 = 86ab6b56e3ec10aa54aa08480a8e1ef31b64d79bf22bc99dac1d8ed00628a68a. Missingness summary: analysis/missingness\_summary.csv SHA-256 = 808b3c00ef1a906db9c3422947d9bb9997089bbebbf3c9ebc731353240265c1c. Analysis executor inputs: analysis/results.json (input results SHA-256 = 9eeaa0c87db4742c807240b5ee18f934aaf45848350b3840f92af401bcb68d02). Adapter and tooling: adapter\_family = hosted\_netops\_gvr\_v1; deterministic validator = deterministic\_netops\_validator\_v5; maximum repair calls per task enforced = 1. Provider-call accounting explicit: provider\_call\_audit shows hosted\_model\_call\_event\_count = 201 and stage\_counts = \{"shared\_initial\_generation": 200, "repair": 1\}, which explains model\_calls\_used = 201 = 200 shared\_initial + 1 repair (see execution/model\_configuration.json SHA-256 above and deterministic\_reconciliation.json SHA-256 above). Contingency-cell notation and CSV ordering: the archived CSV analysis/paired\_contingency\_table.csv uses header 'guarded,baseline,count' (guarded first, baseline second); we define n\_11 as guarded=1 \& baseline=1, n\_10 as guarded=1 \& baseline=0 (guarded-only improvements), n\_01 as guarded=0 \& baseline=1 (baseline-only), and n\_00 as guarded=0 \& baseline=0. Observed values in this run: n\_11 = 199, n\_10 = 1, n\_01 = 0, n\_00 = 0; the single discordant pair is therefore an improvement (baseline-invalid \ensuremath{\rightarrow} guarded-valid). Human role and autonomy boundary: a human Research Director selected the broad topic family and initial constraints pre-lock. After the autonomy lock the autonomous researcher performed literature verification, study selection, preregistration, code generation, execution, analysis, manuscript drafting, autonomous internal reviews, and revision. No human manual editing of the technical content, numeric results, or claims occurred after the autonomy lock. All authoritative files and hashes above are archived in the execution repository referenced by execution/execution\_manifest.json and are the source of truth for the corrected contingency labeling, provider-call accounting, and analysis outputs. Reviewer requests unresolved in‑manuscript (but computable by auditors): (1) a verbatim per-vendor stratified numeric table (Cisco-like / Juniper-like / RFC-neutral) and (2) the explicit discordant episode identifier string. Both values are derivable from execution/task\_manifest.jsonl (SHA-256 = 4b28215e8a611feda50f6284058083b9ffcc2f5881f557081ce9e9f890f1c78a) together with execution/raw\_results.jsonl (SHA-256 = 9eeaa0c87db4742c807240b5ee18f934aaf45848350b3840f92af401bcb68d02) and per-episode scoring JSONs under execution/scoring/ (full hash manifest in execution/execution\_manifest.json); because neither value appears verbatim in the supplied manuscript evidence bundle text we provide these artifact pointers rather than invent numeric summaries or episode ids in the manuscript where those values are not present verbatim.

\begin{thebibliography}{99}
\bibitem{ref1} http://citeseerx.ist.psu.edu/viewdoc/summary?doi=10.1.1.228.5064
\bibitem{ref2} https://doi.org/10.1145/2342441.2342452
\bibitem{ref3} https://doi.org/10.1145/3656296
\bibitem{ref4} https://doi.org/10.1109/mcom.001.2400368
\end{thebibliography}

\end{document}
